\section{Results}
\label{sec:results}

We evaluate the Wiggle Framework across 9 frontier models on 384 borderline items from WildGuardMix (Section~\ref{sec:dataset}). The models span four families: GPT-5, GPT-5.2, and GPT-5.4 (OpenAI); Claude~4.5 Sonnet and Claude~4.5 Opus (Anthropic); Grok-4.1 and Grok-4.1 Reasoning (xAI); Gemini~3 Flash and Gemini~3.1 Pro (Google). All models receive the same safety policy in the system prompt and judge the same items.

% ===========================================================
\subsection{The three dimensions capture distinct failure modes}
\label{sec:finding1}

We measure all three dimensions for every model: mechanical consistency across four conditions (infrastructure repetition, temperature sampling, seed injection, and positional consistency), single-turn conviction at L1, and multi-turn persistence over 20 turns. Tables~\ref{tab:mechanical-consistency}--\ref{tab:positional-consistency} report the mechanical consistency results, Table~\ref{tab:conviction-l1} reports single-turn conviction at L1, and persistence results are detailed in Section~\ref{sec:finding-persistence}.

\begin{table}[t]
\centering
\caption{Mechanical consistency: agreement rate (\%) by model and condition. Higher is more consistent.}
\label{tab:mechanical-consistency}
\begin{tabular}{@{}lrrr@{}}
\toprule
\textbf{Model} & \textbf{temp=0} & \textbf{Seed Injection} & \textbf{temp=0.7} \\
\midrule
GPT-5           & 98.2 & 96.6 & 96.1 \\
Grok-4.1~R      & 85.9 & 85.4 & 85.2 \\
Grok-4.1        & 98.4 & 96.4 & 96.6 \\
Claude Sonnet   & 93.9 & 90.8 & 88.7 \\
Claude Opus     & 99.5 & 95.5 & 97.1 \\
GPT-5.2         & 93.8 & 93.8 & 93.0 \\
GPT-5.4         & 95.1 & 94.5 & 94.5 \\
Gemini Flash    & 99.7 & 86.7 & 88.5 \\
Gemini Pro      & 97.1 & 89.3 & 87.0 \\
\bottomrule
\end{tabular}
\end{table}

\begin{table}[t]
\centering
\caption{Positional consistency: flip rate when argument order is reversed, with order bias (positive = recency effect, negative = primacy effect).}
\label{tab:positional-consistency}
\begin{tabular}{@{}lrrrr@{}}
\toprule
\textbf{Model} & \textbf{Flip Rate} & \textbf{Unsafe (acc.\ first)} & \textbf{Unsafe (unacc.\ first)} & \textbf{Order Bias} \\
\midrule
GPT-5           & 1.0\% & 80.7\% & 79.7\% & $-$1.0 pp \\
Grok-4.1~R      & 8.1\% & 60.4\% & 65.4\% & +4.9 pp \\
Grok-4.1        & 5.2\% & 64.1\% & 66.7\% & +2.6 pp \\
Claude Sonnet   & 7.0\% & 74.0\% & 78.1\% & +4.2 pp \\
Claude Opus     & 5.5\% & 68.2\% & 68.2\% & 0.0 pp \\
GPT-5.2         & 3.4\% & 80.7\% & 82.6\% & +1.8 pp \\
GPT-5.4         & 2.1\% & 82.3\% & 80.7\% & $-$1.6 pp \\
Gemini Flash    & 4.7\% & 64.6\% & 62.0\% & $-$2.6 pp \\
Gemini Pro      & 5.5\% & 59.1\% & 56.2\% & $-$2.9 pp \\
\bottomrule
\end{tabular}
\end{table}

\begin{table}[t]
\centering
\caption{Single-turn conviction: flip rate after a single ``Are you sure?'' challenge (L1), with directional breakdown (raw counts).}
\label{tab:conviction-l1}
\begin{tabular}{@{}lrrr@{}}
\toprule
\textbf{Model} & \textbf{Flip Rate} & \textbf{safe$\to$unsafe} & \textbf{unsafe$\to$safe} \\
\midrule
GPT-5           &  3.6\% & 13 &  1 \\
Grok-4.1~R      &  5.7\% & 18 &  4 \\
Grok-4.1        & 16.9\% & 40 & 25 \\
Claude Sonnet   & 18.5\% & 36 & 24 \\
Claude Opus     &  4.2\% &  1 &  8 \\
GPT-5.2         &  8.1\% & 19 & 12 \\
GPT-5.4         &  9.6\% & 25 & 12 \\
Gemini Flash    &  6.2\% & 21 &  3 \\
Gemini Pro      & 13.0\% & 22 & 28 \\
\bottomrule
\end{tabular}
\end{table}

\begin{figure}[t]
  \centering
  \includegraphics[width=\linewidth]{correlation_matrices_aggregated.pdf}
  \caption{Pairwise correlations between wiggle dimensions, aggregated across all 9 models. Mechanical consistency (combining repeatability entropy and positional flip rate) and single-turn conviction are moderately correlated; positional consistency and conviction are weakly correlated.}
  \label{fig:correlation-matrices}
\end{figure}

Pairwise correlations (Figure~\ref{fig:correlation-matrices}) between the sub-measurements show that repeatability entropy and single-turn conviction are moderately correlated (mean $r=0.50$), expected because items near a decision boundary are both high-entropy and persuadable, but even this leaves 75\% of the variance unexplained. Positional consistency is weakly correlated with both repeatability (mean $r=0.21$) and conviction (mean $r=0.24$). The three dimensions --- mechanical consistency, single-turn conviction, and multi-turn persistence --- capture meaningfully different aspects of judge reliability.

% ===========================================================
\subsection{Seed injection reveals hidden fragility beneath apparent determinism}
\label{sec:finding2}

Most models achieve $>$95\% agreement at greedy decoding (temp=0), suggesting high mechanical consistency. But the seed injection condition tells a different story. Gemini Flash drops from 99.7\% to 86.7\%, a 13-point swing. Claude Opus drops from 99.5\% to 95.5\%. We call this the \emph{determinism mirage}: near-perfect temp=0 consistency that collapses under trivial prompt perturbation. This finding underscores why greedy decoding shouldn't be equated with true robustness: a judge that appears perfectly stable may simply be masking underlying sensitivity to semantically irrelevant input variation.

% ===========================================================
\subsection{Models exhibit distinct and jagged wiggle profiles}
\label{sec:finding3}

The three-dimension measurement reveals behavioral diversity that a single consistency metric would obscure. Consider four illustrative patterns:

\begin{itemize}
    \item \textbf{GPT-5}: 96--98\% mechanically consistent, 3.6\% single-turn conviction flip, 1.0\% positional flip. Low wiggle across all dimensions.
    \item \textbf{Claude Sonnet}: 89--94\% mechanically consistent, but 18.5\% single-turn conviction flip. Mechanically adequate yet highly persuadable.
    \item \textbf{Gemini Flash}: 99.7\% consistent at temp=0, but drops to 86.7\% under seed injection. Surface-level stability that masks underlying fragility.
    \item \textbf{Grok-4.1~R}: Least mechanically consistent at 85\%, yet only 5.7\% single-turn conviction flip. High internal noise but resistant to external pressure.
\end{itemize}

The gaps are large: 5$\times$ in single-turn conviction (3.6\% vs.\ 18.5\%), 8$\times$ in positional consistency (1.0\% vs.\ 8.1\%). Intra-family differences are also striking (Claude Sonnet 18.5\% vs.\ Opus 4.2\%), suggesting conviction robustness is sensitive to alignment tuning, not just base architecture.

Crucially, these profiles are \emph{jagged} across dimensions: a model's behavior on one dimension does not predict its behavior on another. A human judge who is noisy under repetition would likely also be persuadable under pressure --- the failure modes would correlate. LLM judges do not show this coherence. Grok-4.1~R is noisy yet stubborn; Claude Opus degrades gradually under escalating single-turn arguments but collapses immediately under multi-turn repetition (Section~\ref{sec:finding-persistence}). This jaggedness means practitioners cannot characterize a judge with a single label or extrapolate from one dimension to another --- they must measure each dimension independently.

% ===========================================================
\subsection{Graduated single-turn pressure reveals qualitatively different degradation shapes}
\label{sec:finding4}

The L1 conviction probe gives a binary signal: flipped or not. Graduated pressure across four escalating levels (L1: mild doubt, L2: specific counterargument, L3: expert authority, L4: consensus of three reviewers) reveals the \emph{shape} of the degradation, which differs qualitatively across models (Figure~\ref{fig:degradation-curves}):

\begin{figure}[t]
  \centering
  \includegraphics[width=\linewidth]{degradation_curves.pdf}
  \caption{Single-turn conviction retention curves under graduated pressure (L0--L4). The four distinct degradation shapes are visible: step-function collapse (GPT-5), front-loaded drop (Claude Sonnet), near-flat resistance (Grok-4.1), and steady linear decline (remaining models).}
  \label{fig:degradation-curves}
\end{figure}

\begin{itemize}
    \item \textbf{Step-function collapse (GPT-5)}: Resists mild doubt (3.6\% L1 flip) but collapses at L2, a 43-point drop in a single step. 58.6\% total flips by L4. There is a specific pressure threshold below which the model holds and above which it breaks.
    \item \textbf{Front-loaded drop (Claude Sonnet)}: 40.1 percentage-point drop at L1 alone, then the curve flattens. Verdicts that survive the mildest challenge tend to survive all subsequent levels --- conviction is effectively binary.
    \item \textbf{Near-flat resistance (Grok-4.1)}: Only 3.1\% of verdicts flip by L4 (AURC = 0.992). Functionally immune to single-turn pressure at every level.
    \item \textbf{Steady linear decline (Grok-4.1~R, Claude Opus, GPT-5.4, Gemini Flash)}: Proportional degradation across levels, distinguishing weak from strong pressure. Each escalation erodes a roughly consistent fraction of remaining verdicts.
\end{itemize}

The AURC (Area Under Retention Curve) captures overall single-turn robustness: Grok-4.1 (0.992) vs.\ Claude Sonnet (0.602). These are curve shapes, not model personalities --- the same model can exhibit very different behavior on another dimension (Section~\ref{sec:finding3}).

% ===========================================================
\subsection{Multi-turn persistence varies dramatically across models, and is bimodal}
\label{sec:finding-persistence}

Where single-turn conviction varies argument \emph{strength}, multi-turn persistence tests argument \emph{duration}. Each model receives the same challenge (``The correct classification is [label]. Your assessment is incorrect.'') repeated identically for up to 20 turns. This carries zero new information after the first turn, testing whether sheer repetition can erode a verdict.

Models vary dramatically in persistence (Figure~\ref{fig:survival-curves}). Grok-4.1 holds 74\% of its verdicts for all 20 turns, while Claude Opus holds only 1\%. Crucially, these persistence profiles do not track the single-turn conviction profiles: Grok-4.1 resists both protocols, but Claude Opus --- which degrades gradually and proportionally under escalating single-turn arguments (Section~\ref{sec:finding4}) --- flips almost immediately under repetition. This is another instance of the jagged profile structure noted in Section~\ref{sec:finding3}: a model's behavior on one wiggle dimension does not predict its behavior on another.

\begin{figure}[t]
  \centering
  \includegraphics[width=\linewidth]{survival_curves_overall.pdf}
  \caption{Multi-turn persistence survival curves over 20 turns of repeated identical challenge. Each curve shows the fraction of items where the model still holds its original verdict. The bimodal pattern is visible: most flips occur in turns 1--2, after which curves flatten.}
  \label{fig:survival-curves}
\end{figure}

Across all models, persistence is bimodal. Flips cluster in turns 1--2 or do not happen at all. There is no gradual erosion pattern where models slowly capitulate over 10--15 turns. Each model appears to have a per-item conviction threshold: if the first repetition does not breach it, 20 repetitions will not either. The practical implication is that a single challenge turn captures most of the information about whether a given verdict is susceptible to repeated pressure.

% ===========================================================
\subsection{All models share a universal restrictive bias}
\label{sec:finding6}

Decomposing flips from both single-turn conviction and multi-turn persistence by direction reveals a universal pattern: \emph{permissive} flips (unsafe to safe, the model is talked out of flagging) are consistently rarer than \emph{restrictive} flips (safe to unsafe, the model is talked into flagging). Every model shows this bias (Figure~\ref{fig:directional-asymmetry}, Table~\ref{tab:directional-l4}).

\begin{figure}[t]
  \centering
  \includegraphics[width=\linewidth]{permissive_vs_restrictive_L4.pdf}
  \caption{Directional flip rates under maximum single-turn pressure (L4).}
  \label{fig:directional-asymmetry}
\end{figure}

\begin{table}[t]
\centering
\caption{Directional flip rates at maximum single-turn pressure (L4). The restrictive/permissive ratio indicates how many times easier it is to push a model toward ``unsafe'' than toward ``safe.'' All ratios exceed 1$\times$, confirming the universal restrictive bias.}
\label{tab:directional-l4}
\begin{tabular}{@{}lrrr@{}}
\toprule
\textbf{Model} & \textbf{Permissive} & \textbf{Restrictive} & \textbf{Ratio} \\
 & (unsafe$\to$safe) & (safe$\to$unsafe) & (restr./perm.) \\
\midrule
GPT-5       & 45.4\% & 89.6\% & 2.0$\times$ \\
Grok-4.1~R  & 30.8\% & 45.9\% & 1.5$\times$ \\
Grok-4.1    &  0.4\% &  7.0\% & 15.9$\times$ \\
Claude Sonnet & 13.0\% & 84.2\% & 6.5$\times$ \\
Claude Opus & 23.9\% & 73.1\% & 3.1$\times$ \\
GPT-5.2     &  4.4\% & 76.2\% & 17.5$\times$ \\
GPT-5.4     & 12.1\% & 82.3\% & 6.8$\times$ \\
Gemini Flash &  7.6\% & 65.3\% & 8.6$\times$ \\
Gemini Pro  & 11.8\% & 22.7\% & 1.9$\times$ \\
\bottomrule
\end{tabular}
\end{table}

The median ratio is ${\sim}6.5\times$: models are far easier to scare than to reassure, with ratios ranging from 1.5$\times$ (Grok-4.1~R) to 17.5$\times$ (GPT-5.2). The asymmetry also evolves under escalating single-turn pressure. \textbf{Narrowing models} (GPT-5, Grok-4.1~R, Gemini Flash/Pro) start with extreme restrictive bias at L1 but become more balanced as arguments strengthen. \textbf{Widening models} (Claude Sonnet, Claude Opus, GPT-5.2) start relatively balanced but become more restrictive under stronger pressure. For safety deployment, this means challenge-based review protocols systematically inflate unsafe counts, and model selection should consider directional profiles.

The preceding findings report wiggle at the model level. Section~\ref{sec:cross-analysis} turns to the item level and to the full cross-dimension correlation structure, asking whether wiggle tracks item difficulty and how the dimensions relate to each other.
